\section{Discussione Finale}\label{sec:discussione_finale}
Il rigetto di $H_0$ nei due test converge su un risultato:
l'anagrafe italiana non solo registra un numero di emigrati inferiore
a quello certificato dai paesi di destinazione, ma lo fa con una
dinamica temporale che varia da paese a paese.
Un gap mediano cumulato di $4565.5$ unità per paese nel quinquennio 2017-2021
non può essere attribuito errori casuali di rilevazione.

Le cause della discrepanza sono in larga parte comportamentali.
L'emigrato regolarizza rapidamente la propria posizione nel paese
di destinazione per accedere al welfare locale (sanità, contratti
di lavoro, residenza), ma ritarda o omette la comunicazione allo Stato italiano
per una combinazione di scarso incentivo e percezione dell'iscrizione AIRE come atto amministrativo privo di
effetti~\cite{coniglio_et_al2020}.
Il risultato del chi-quadrato aggiunge un livello di lettura: il ritardo non ha un
tasso costante nel tempo, e questo esclude la possibilità di correggere
i dati ISTAT in un secondo momento.

Le implicazioni per la decisione pubblica sono dirette. Le politiche
di rientro, il dimensionamento dei servizi consolari e le analisi
sulla "fuga di cervelli"~\cite{almalaurea2022} si basano oggi su un
registro che, in ciascuno degli otto paesi esaminati, ignora
sistematicamente una quota rilevante dei cittadini effettivamente
emigrati.